\chapter{DEPLOYMENT AND OPERATIONS}

\noindent\textbf{Mục tiêu chương}

\begin{enumerate}[label=-]
\item Tách phần triển khai và vận hành khỏi chapter output/release để operator tìm đúng chỗ nhanh hơn.
\item Chốt baseline host mới, runtime configuration, manual smoke, và schedule enable flow.
\item Đưa guidance của technical PDF tiến gần hơn kiểu operator docs, nhưng vẫn giữ được tính handover-ready.
\item Ghi rõ các guardrail vận hành để giảm lỗi cấu hình khi triển khai sang máy mới.
\end{enumerate}

\section{Vai trò của chapter deployment and operations}

Chapter này trả lời nhóm câu hỏi thực tế khi đem project sang một host Ubuntu / \texttt{systemd} mới:

\begin{enumerate}[label=-]
\item clone và bootstrap theo baseline nào;
\item các giá trị runtime nào bắt buộc phải thay theo từng hệ thống;
\item nên manual smoke ra sao trước khi bật scheduler;
\item cài và xác nhận \texttt{systemd} service / timer thế nào để vận hành ổn định.
\end{enumerate}

Nội dung chương này phải được giữ đồng bộ với \texttt{docs/workflows/deployment\_runbook.md} vì đó là nguồn operator-facing chi tiết nhất cho thao tác triển khai thật.

\section{New-host bootstrap baseline}

Baseline triển khai đang được khuyến nghị cho dự án là:

\begin{enumerate}[label=-]
\item host Ubuntu có \texttt{systemd};
\item service account \texttt{energy-report};
\item project root \texttt{/srv/energy-report};
\item một \texttt{venv} mới được tạo trực tiếp trên host;
\item Chrome hoặc Chromium hiện diện để PDF export chạy được;
\item output path writable và không nằm dưới hidden path.
\end{enumerate}

Các guardrail bootstrap cần giữ:

\begin{enumerate}[label=-]
\item không copy \texttt{venv/} từ máy cũ sang máy mới;
\item chỉ dùng sample \texttt{systemd} units trực tiếp nếu host bám đúng baseline user/path;
\item nếu host phải chạy theo giờ Việt Nam, cần xác nhận host timezone trước khi bật timer;
\item giữ rõ ranh giới giữa canonical output path và staging path dùng cho PDF export.
\end{enumerate}

\section{Operator-first runtime configuration}

\subsection{Host-specific values cần rà soát}

Để tránh việc operator phải mở \texttt{config/.env} rồi sửa thủ công từng dòng, deployment runbook hiện đã có block terminal paste-ready. Trước khi paste, người cài chỉ cần xác nhận các giá trị host-specific sau:

\begin{enumerate}[label=-]
\item \texttt{MYSQL\_HOST};
\item \texttt{MYSQL\_PORT};
\item \texttt{MYSQL\_DATABASE};
\item \texttt{MYSQL\_USER};
\item \texttt{MYSQL\_PASSWORD};
\item \texttt{OUTPUT\_DIR} nếu host không dùng \texttt{/srv/energy-report-output};
\item \texttt{PRINT\_STAGING\_DIR} nếu staging cần khác output root;
\item khi cần, \texttt{WORKSHOP\_NAME} và \texttt{ENERGY\_UNIT}.
\end{enumerate}

\subsection{Paste-ready block để ghi \texttt{config/.env}}

Block bên dưới phản ánh operator-first flow hiện tại. Canonical source vẫn là \texttt{docs/workflows/deployment\_runbook.md}, nhưng technical PDF giữ lại block này để người vận hành đọc tài liệu bàn giao vẫn nhìn được đúng mẫu cấu hình cần paste.

\begingroup
\small
\begin{verbatim}
sudo tee /srv/energy-report/config/.env >/dev/null <<'EOF'
# --- Change these values to match the target system ---
MYSQL_HOST=__FILL_DB_HOST__
MYSQL_PORT=3306
MYSQL_DATABASE=__FILL_DB_NAME__
MYSQL_USER=__FILL_DB_USER__
MYSQL_PASSWORD=__FILL_DB_PASSWORD__

# --- Change these paths only if your host uses a different storage layout ---
OUTPUT_DIR=/srv/energy-report-output
PRINT_STAGING_DIR=/srv/energy-report-output

# --- Usually keep these values as-is ---
REPORT_FILENAME=energy_automatic_report
LOG_LEVEL=INFO
WORKSHOP_NAME=ENERGY REPORT
ENERGY_UNIT=kWh

# --- Keep blank in normal scheduled production mode ---
REPORT_ANCHOR_DATE=
EOF

sudo chown energy-report:energy-report /srv/energy-report/config/.env
sudo chmod 640 /srv/energy-report/config/.env
sudo mkdir -p /srv/energy-report/logs /srv/energy-report-output
sudo chown -R energy-report:energy-report \
  /srv/energy-report/logs \
  /srv/energy-report-output
\end{verbatim}
\endgroup

Các điểm bắt buộc phải nhớ khi dùng block này:

\begin{enumerate}[label=-]
\item \texttt{REPORT\_ANCHOR\_DATE} phải để trống trong steady-state scheduled mode;
\item nếu giữ nhầm một anchor date cố định sau smoke test, timer sẽ lặp lại ngày đó thay vì current date thực;
\item \texttt{OUTPUT\_DIR} và \texttt{PRINT\_STAGING\_DIR} phải là non-hidden writable paths;
\item nếu host không dùng baseline \texttt{/srv/energy-report}, block và \texttt{systemd} units phải được chỉnh đồng bộ.
\end{enumerate}

\section{Manual smoke trước khi bật schedule}

Trước khi enable timer, host mới nên đi qua ít nhất một vòng manual smoke:

\begin{enumerate}[label=-]
\item chạy \texttt{./venv/bin/python -m src.main};
\item xác nhận daily artifacts được tạo ra đúng canonical month folder;
\item khi cần, dùng các anchor đã chấp nhận là \texttt{2025-06-25}, \texttt{2025-06-29}, và \texttt{2025-05-31} để kiểm tra daily-only, Sunday, và month-end behavior;
\item sau khi smoke xong, trả \texttt{REPORT\_ANCHOR\_DATE} về trạng thái trống trước khi chuyển sang schedule thật.
\end{enumerate}

\section{Systemd installation và fixed daily schedule}

Runtime hiện tại được vận hành an toàn nhất bằng \texttt{systemd} \textit{oneshot service} kết hợp \texttt{timer}. Repo đang cung cấp ba artifact hỗ trợ flow này:

\begin{enumerate}[label=-]
\item \texttt{deploy/systemd/energy-report-etl.service};
\item \texttt{deploy/systemd/energy-report-etl.timer};
\item \texttt{deploy/systemd/install\_systemd\_units.sh}.
\end{enumerate}

Helper script dùng để giảm thao tác tay khi host bám đúng baseline. Nó tự backup unit cũ nếu có, copy sample files sang \texttt{/etc/systemd/system/}, reload \texttt{systemd}, rồi enable timer trong một lần chạy.

Baseline service/timer hiện hành gồm:

\begin{enumerate}[label=-]
\item \texttt{User=energy-report} và \texttt{Group=energy-report};
\item \texttt{WorkingDirectory=/srv/energy-report};
\item \texttt{ExecStart=/srv/energy-report/venv/bin/python -m src.main};
\item \texttt{OnCalendar=*-*-* 23:00:00};
\item \texttt{Persistent=true}.
\end{enumerate}

Lệnh xác nhận tối thiểu sau khi cài:

\begin{enumerate}[label=-]
\item \texttt{sudo systemctl status energy-report-etl.service --no-pager};
\item \texttt{sudo systemctl status energy-report-etl.timer --no-pager};
\item \texttt{systemctl list-timers energy-report-etl.timer --all};
\item \texttt{journalctl -u energy-report-etl.service -n 100 --no-pager}.
\end{enumerate}

\section{Operational validation checklist}

Trước khi coi host mới là usable cho production-like schedule, nên xác nhận tối thiểu:

\begin{enumerate}[label=-]
\item MySQL credentials đúng và truy cập được dưới service account;
\item browser runtime có sẵn;
\item output path writable;
\item manual run thành công;
\item service start thủ công thành công;
\item timer đã enabled và hiển thị next run đúng 23:00 theo timezone mong muốn;
\item operator biết rõ block paste-ready nào là canonical source cần dùng lại khi cài host tiếp theo.
\end{enumerate}

\section{Tổng kết chương}

Chapter này tách riêng phần triển khai và vận hành để technical manual gần hơn với operator docs. Sau chapter output/release, người đọc giờ có một chỗ độc lập để tra baseline máy mới, runtime configuration, manual smoke flow, và \texttt{systemd} schedule setup mà không phải lọc chúng ra khỏi troubleshooting content.
